\clearpage
\VerdierPhysicalPageAnchor{0270}%
\begingroup
\small
\setlength{\parindent}{0pt}%
\setlength{\parskip}{0pt}%
\setlength{\emergencystretch}{3em}%
\newcommand{\VerdierContentsRow}[3]{%
  \par\noindent
  \hangindent=#1\relax\hangafter=1
  \hspace*{#1}#2\nobreak
  \leaders\hbox{\kern.32em.\kern.32em}\hfill
  \nobreak\hbox{\kern.65em #3}\par}
\newcommand{\VerdierContentsWrappedRow}[3]{%
  \par\noindent\hspace*{.5em}#1\par\nobreak
  \noindent\hspace*{3.15em}#2\nobreak
  \leaders\hbox{\kern.32em.\kern.32em}\hfill
  \nobreak\hbox{\kern.65em #3}\par}
\newcommand{\VerdierContentsSection}[2]{%
  \par\addvspace{.55\baselineskip}\noindent
  \hfill\textbf{#1.}\enspace #2\hfill\mbox{}\par\smallskip}

\VerdierSemanticAnchor{toc-ch3-sec3-2}%
\VerdierContentsRow{.5em}{\textbf{3.2.}\enspace La construction de Yoneda}{193}
\VerdierSemanticAnchor{toc-ch3-sec3-3}%
\VerdierContentsRow{.5em}{\textbf{3.3.}\enspace La catégorie $\mathcal A^{\operatorname{Ext}}$}{200}

\VerdierSemanticAnchor{toc-ch3-sec4}%
\VerdierContentsSection{4}{Objets spectraux usuels.}
\VerdierSemanticAnchor{toc-ch3-sec4-1}%
\VerdierContentsRow{.5em}{\textbf{4.1.}\enspace Premier objet spectral canonique}{207}
\VerdierSemanticAnchor{toc-ch3-sec4-2}%
\VerdierContentsWrappedRow{\textbf{4.2.}\enspace Objet spectral à valeurs dans $D(\mathcal A)$ associé à un objet}{filtré de $\operatorname{comp}(\mathcal A)$}{217}
\VerdierSemanticAnchor{toc-ch3-sec4-3}%
\VerdierContentsRow{.5em}{\textbf{4.3.}\enspace Deuxième objet spectral canonique}{219}
\VerdierSemanticAnchor{toc-ch3-sec4-4}%
\VerdierContentsRow{.5em}{\textbf{4.4.}\enspace Suites spectrales usuelles}{221}
\VerdierSemanticAnchor{toc-ch3-sec4-5}%
\VerdierContentsRow{.5em}{\textbf{4.5.}\enspace Problèmes de convergence}{223}
\VerdierSemanticAnchor{toc-ch3-sec4-6}%
\VerdierContentsRow{.5em}{\textbf{4.6.}\enspace Objet spectral de Cartan--Eilenberg}{227}

\vspace{1.2\baselineskip}
\VerdierSemanticAnchor{toc-index-notation}%
\VerdierContentsRow{0pt}{\textbf{Index des notations}}{235}
\vspace{.8\baselineskip}
\VerdierSemanticAnchor{toc-index-terminological}%
\VerdierContentsRow{0pt}{\textbf{Index terminologique}}{241}
\vspace{.8\baselineskip}
\VerdierSemanticAnchor{toc-bibliography}%
\VerdierContentsRow{0pt}{\textbf{Bibliographie}}{249}
\vspace{.8\baselineskip}
\VerdierSemanticAnchor{toc-contents}%
\VerdierContentsRow{0pt}{\textbf{Table des matières}}{251}
\endgroup
